% Generated from bibliography.tex by build.py.
\begin{thebibliography}{99}
\addcontentsline{toc}{section}{References}

\bibitem{bkn2014}
A.~D.~Barbour, E.~Kowalski, and A.~Nikeghbali,
\emph{Mod-discrete expansions},
Probability Theory and Related Fields \textbf{158}
(2014), 859--893.
\href{https://doi.org/10.1007/s00440-013-0498-8}
{doi:10.1007/s00440-013-0498-8}.
Section numbers refer to the 26-page author text, Sections~5 and~8.1:
\url{https://people.math.ethz.ch/~kowalski/mod-discrete.pdf}.

\bibitem{binyaminicluckersnovikov2024}
G.~Binyamini, R.~Cluckers, and D.~Novikov,
\emph{Bounds for rational points on algebraic curves, optimal
in the degree, and dimension growth},
International Mathematics Research Notices \textbf{2024} (2024),
no.~11, 9256--9265.
\href{https://doi.org/10.1093/imrn/rnae034}
{doi:10.1093/imrn/rnae034}.
The degree-uniform integer-point estimate is cited from
\href{https://arxiv.org/abs/2302.04209v2}
{arXiv:2302.04209v2} (20 September 2023),
Theorem~2, equation~(27).

\bibitem{bretwngxu2026}
R.~de~la~Bret\`eche, V.~Y.~Wang, and M.~W.~Xu,
\emph{Random multiplicative functions and making squares from
polynomial values},
\href{https://arxiv.org/abs/2607.06398v1}
{arXiv:2607.06398v1} (2026), Theorem~1.2 and
Lemmas~6.2--6.3.

\bibitem{erdosturk1984}
P.~Erd\H{o}s and J.~Turk,
\emph{Products of integers in short intervals},
Acta Arithmetica \textbf{44} (1984), no.~2, 147--174.
\href{https://doi.org/10.4064/aa-44-2-147-174}
{doi:10.4064/aa-44-2-147-174}.
\href{https://www.renyi.hu/~p_erdos/1984-17.pdf}
{Article text}.

\bibitem{heathbrown2012}
D.~R.~Heath-Brown,
\emph{Square-free values of \(n^2+1\)},
Acta Arithmetica \textbf{155} (2012), no.~1, 1--13.
\href{https://doi.org/10.4064/aa155-1-1}
{doi:10.4064/aa155-1-1}.
Preprint version:
\href{https://arxiv.org/abs/1010.6217v3}
{arXiv:1010.6217v3}, Lemma~2.

\bibitem{munsch2026}
M.~Munsch, I.~E.~Shparlinski, Y.-C.~Sun, and Y.~Xiao,
\emph{A large sieve inequality for sums of Legendre symbols
over short intervals},
\href{https://arxiv.org/abs/2604.23661v3}
{arXiv:2604.23661v3} (5 August 2026), Lemma~3.3 and its proof.

\bibitem{pila1995}
J.~Pila,
\emph{Density of integral and rational points on varieties},
in \emph{Columbia University Number Theory Seminar
(New York, 1992)}, Ast\'erisque \textbf{228} (1995),
183--187.
\href{https://www.numdam.org/item/AST_1995__228__183_0/}
{Numdam: AST\_1995\_\_228\_\_183\_0}.

\bibitem{reuss2014}
T.~Reuss,
\emph{Pairs of \(k\)-free numbers, consecutive square-full
numbers},
\href{https://arxiv.org/abs/1212.3150v2}
{arXiv:1212.3150v2} (19 March 2014), Sections~2.1--2.3.

\bibitem{salberger2023}
P.~Salberger,
\emph{Counting integral points of affine hypersurfaces},
\href{https://arxiv.org/abs/2311.05561v1}
{arXiv:2311.05561v1} (9 November 2023), Lemma~2.

\bibitem{zacharovashwang2010}
V.~Zacharovas and H.-K.~Hwang,
\emph{A Charlier--Parseval approach to Poisson approximation
and its applications},
Lithuanian Mathematical Journal \textbf{50} (2010), no.~1,
88--119.
\href{https://doi.org/10.1007/s10986-010-9073-5}
{doi:10.1007/s10986-010-9073-5}.
Preprint version used for equation and remark numbers:
\href{https://arxiv.org/abs/0810.4756v1}
{arXiv:0810.4756v1}, (2.13) and Remark~4.4.

\bibitem{griesmer1960}
J.~H.~Griesmer,
\emph{A bound for error-correcting codes},
IBM Journal of Research and Development \textbf{4} (1960),
no.~5, 532--542.
\href{https://doi.org/10.1147/rd.45.0532}
{doi:10.1147/rd.45.0532}.
\href{https://bitsavers.trailing-edge.com/pdf/ibm/IBM_Journal_of_Research_and_Development/045/ibmrd0405M.pdf}
{Original article scan}, Section~2, equation~(2.10).

\bibitem{cilleruelogaraev2011}
J.~Cilleruelo and M.~Z.~Garaev,
\emph{Concentration of points on two and three dimensional
modular hyperbolas and applications},
Geometric and Functional Analysis \textbf{21} (2011),
no.~4, 892--904.
\href{https://doi.org/10.1007/s00039-011-0127-6}
{doi:10.1007/s00039-011-0127-6}.
Numbering refers to
\href{https://arxiv.org/abs/1007.1526v2}{arXiv:1007.1526v2},
Section~3, Lemma~4 and Proposition~1.

\bibitem{tao2026anatomy}
T.~Tao,
\emph{Products of consecutive integers with unusual anatomy},
\href{https://arxiv.org/abs/2603.27990v3}{arXiv:2603.27990v3}
(26 September 2026), Lemma~2.10.

\bibitem{shparlinskisleiman2026}
I.~E.~Shparlinski and N.~Sleiman,
\emph{Counting consecutive multiplicatively dependent triples},
\href{https://arxiv.org/abs/2609.29408v1}
{arXiv:2609.29408v1} (24 September 2026),
Theorem~1.1, Lemmas~2.1--2.2, and Section~3.

\bibitem{binyaminicluckerskato2025}
G.~Binyamini, R.~Cluckers, and F.~Kato,
\emph{Sharp bounds for the number of rational points on algebraic curves
and dimension growth, over all global fields},
Proceedings of the London Mathematical Society \textbf{130} (2025),
no.~1, e70016. Published online 16 December 2024.
\href{https://doi.org/10.1112/plms.70016}
{doi:10.1112/plms.70016}.
Corollary~8, equation~(23).

\bibitem{bretkurlbergshparlinski2021}
R.~de~la~Bret\`eche, P.~Kurlberg, and I.~E.~Shparlinski,
\emph{On the number of products which form perfect powers and
discriminants of multiquadratic extensions},
International Mathematics Research Notices \textbf{2021} (2021),
no.~22, 17140--17169. Published online 30 December 2019.
\href{https://doi.org/10.1093/imrn/rnz316}
{doi:10.1093/imrn/rnz316}.
\href{https://kurlberg.github.io/eprints/SquareProduct.pdf}
{Author text}, Theorem~2.2 and Corollary~2.4.

\end{thebibliography}
